\section{Fourier fields for a smaller unitary}\label{sec:fields}

Let $d=k+s\le 2^n$, with $k,s\ge1$, and suppose that
\begin{equation}\label{fields:block}
 V=\begin{pmatrix}A&B\\ C&D\end{pmatrix}\in U(d),
 \qquad \norm{D}\le r<1,\quad r\ge0.
\end{equation}
The first $k$ coordinates are the ports and the last $s$ coordinates are
internal. We will associate to $V$ a family of $k$-dimensional unitaries
$F(z)$ and two nearly isometric maps from $\C^d$ into the space of
port-valued functions of $z$. These maps intertwine $V$ with the controlled
application of $F(z)$. Their Fourier coefficients will have few nonzero
entries in each row, which makes the maps implementable by slots.

\subsection{Closing the internal coordinates}

For the moment let $Z$ be any diagonal unitary of order $s$. Given a port
input $u\in\C^k$, feed the internal output $v\in\C^s$ back into the
internal input through $Z$. If $y\in\C^k$ is the port output, the equation
$V(u,Zv)=(y,v)$ becomes
\[
 y=Au+BZv,\qquad v=Cu+DZv.
\]
Since $\norm{DZ}<1$, the second equation has the unique solution $v=Xu$,
and the first then gives $y=Fu$, where
\begin{equation}\label{fields:feedback}
 X=(I-DZ)^{-1}C,\qquad F=A+BZX,
 \qquad
 E=\begin{pmatrix}I_k\\ ZX\end{pmatrix},\qquad
 O=\begin{pmatrix}F\\ X\end{pmatrix}.
\end{equation}
Thus $Eu$ and $Ou$ collect the full input and output of $V$ after the
internal coordinates have been closed.
For a scalar phase $Z=zI$, this is the standard transfer function of a
unitary block matrix; see \cite[Section~3]{FKK2009}. The same conservation
of norm proves unitarity with independently chosen diagonal phases.

\begin{lemma}\label{fields:unitary}
For the unitary block matrix in \eqref{fields:block} and any diagonal
unitary $Z$, the matrix $F$ in \eqref{fields:feedback} is unitary.
If
\begin{equation}\label{fields:R}
 R=OF^*=\begin{pmatrix}I_k\\ Y\end{pmatrix},
 \qquad Y=XF^*,
\end{equation}
then
\begin{equation}\label{fields:exact}
 FE^*=R^*V,
 \qquad
 Y=(I-Z^*D^*)^{-1}Z^*B^*.
\end{equation}
\end{lemma}

\begin{proof}
The definitions give $VE=O$. Consequently,
\[
 I_k+X^*X=E^*E=O^*O=F^*F+X^*X,
\]
so $F^*F=I_k$. Because $F$ is square, it is unitary. Thus $OF^*$ has
the stated top block, and $V^*R=EF^*$. Taking adjoints proves the first
identity in \eqref{fields:exact}. Comparing the bottom blocks of
$V^*R=EF^*$ gives
\[
 B^*+D^*Y=ZY.
\]
Multiplication by $Z^*$ and $\norm{Z^*D^*}\le r<1$ give the second identity.
\end{proof}

\subsection{Truncation and averaging}

Fix an integer truncation length $L\ge1$, and set
\begin{equation}\label{fields:parameters}
 b=nL+\lceil\log_2(L+1)\rceil,\qquad
 \mathcal G=(\Z/(2^b))^{L+1}.
\end{equation}
For $w,z\in\mathcal G$, write
\[
 \chi_w(z)=\exp(2\pi i\,w\cdot z/2^b).
\]
Number the internal coordinates by $i=1,\ldots,s$ and put
\begin{equation}\label{fields:phases}
 v_i=(1,i,i^2,\ldots,i^L)\in\mathcal G,
 \qquad Z(z)=\diag\bigl(\chi_{v_1}(z),\ldots,\chi_{v_s}(z)\bigr).
\end{equation}
The constant first coordinate of $v_i$ gives all diagonal entries of $Z(z)$
a common phase when the first coordinate of $z$ varies. Averaging that phase
will separate terms of different lengths in the truncated series. The
remaining coordinates record power sums: the Fourier frequency of a product
of at most $L$ phases will determine its multiset of internal indices, as proved in
Proposition~\ref{fields:sparsity}. These two properties will give the norm
bound and the row bound, respectively.

Apply Lemma~\ref{fields:unitary} at each $z$ to obtain $F(z),E(z),R(z)$.
Truncate the two Neumann expansions by defining
\begin{equation}\label{fields:truncations}
 E_L(z)=\begin{pmatrix}
 I_k\\[2pt] \displaystyle\sum_{\ell=0}^{L-1}Z(DZ)^\ell C
 \end{pmatrix},\qquad
 R_L(z)=\begin{pmatrix}
 I_k\\[2pt] \displaystyle\sum_{\ell=0}^{L-1}(Z^*D^*)^\ell Z^*B^*
 \end{pmatrix}.
\end{equation}
Here and below, products containing $Z$ are evaluated at the same $z$.
Since every block of a unitary has norm at most one,
\begin{equation}\label{fields:tails}
 \norm{E_L(z)-E(z)},\ \norm{R_L(z)-R(z)}
 \le \tau,\qquad \tau=\frac{r^L}{1-r}.
\end{equation}

Define the field maps $\mathcal E_L,\mathcal R_L:
\C^d\longrightarrow\C^{\mathcal G}\otimes\C^k$ by
\begin{equation}\label{fields:maps}
 \mathcal E_L x=\frac1{\sqrt{|\mathcal G|}}
   \sum_{z\in\mathcal G}\ket z\otimes E_L(z)^*x,
 \qquad
 \mathcal R_L x=\frac1{\sqrt{|\mathcal G|}}
   \sum_{z\in\mathcal G}\ket z\otimes R_L(z)^*x.
\end{equation}
Although a single matrix $E_L(z)^*$ need not be a contraction,
averaging its squared norm over $z$ nearly preserves the norm of every input.

\begin{proposition}\label{fields:gram}
For either $\mathcal T=\mathcal E_L$ or $\mathcal T=\mathcal R_L$,
\begin{equation}\label{fields:gram-bound}
 (1-r^{2L})I_d\ \le\ \mathcal T^*\mathcal T\ \le\ I_d.
\end{equation}
In addition, let $\widetilde F$ apply $F(z)$ to the port register,
controlled on $z$. Then
\begin{equation}\label{fields:intertwine}
 \norm{\widetilde F\mathcal E_L-\mathcal R_L V}\le2\tau.
\end{equation}
\end{proposition}

\begin{proof}
The Gram matrix of $\mathcal E_L$ is the average of $E_LE_L^*$.
Fix the last $L$ coordinates of $z$ and vary its first coordinate.
Writing $Z=\omega Z'$, the scalar $\omega$ ranges uniformly over the
$2^b$-th roots of unity, while $Z'$ stays fixed. The summand of index
$\ell$ in the lower block of $E_L$ is
\[
 \omega^{\ell+1}Z'W^\ell C,\qquad W=DZ'.
\]
The degrees $0,1,\ldots,L$ are distinct modulo $2^b$. Their orthogonality
under averaging shows that the port block of $\E_\omega E_LE_L^*$ is
$I_k$, that its off-diagonal blocks vanish, and that its internal block is
\[
 Z'\left(\sum_{\ell=0}^{L-1}W^\ell CC^*(W^*)^\ell\right)(Z')^*.
\]
Unitarity of $V$ gives $CC^*=I-DD^*=I-WW^*$. Therefore the sum telescopes:
\begin{equation}\label{fields:telescope}
 \sum_{\ell=0}^{L-1}W^\ell(I-WW^*)(W^*)^\ell
   =I-W^L(W^*)^L.
\end{equation}
As $\norm{W}\le r$, the internal block lies between
$(1-r^{2L})I_s$ and $I_s$. Averaging the remaining coordinates of $z$
preserves these inequalities.

For $R_L$, set $W=(Z')^*D^*$ and $C_0=(Z')^*B^*$. Its lower summand is
$\omega^{-(\ell+1)}W^\ell C_0$. The negative degrees are again distinct
from each other and from zero. Now
\[
 C_0C_0^*=(Z')^*(I-D^*D)Z'=I-WW^*,
\]
so the internal Gram block is exactly the sum in
\eqref{fields:telescope}, with these choices of $W,C_0$.
This proves \eqref{fields:gram-bound} for both maps.

Finally, Lemma~\ref{fields:unitary} and \eqref{fields:tails} give, pointwise,
\[
 \norm{F(z)E_L(z)^*-R_L(z)^*V}
 \le \norm{E_L(z)-E(z)}+\norm{R_L(z)-R(z)}\le2\tau.
\]
Apply this bound to a fixed vector and average its squared norm over $z$
to obtain \eqref{fields:intertwine}.
\end{proof}

\subsection{Few possible input indices at each Fourier frequency}

The norm bounds supply nearly isometric maps. To implement them, we also
need a restriction on their matrix entries. Use the Fourier basis
\[
 \ket{\widehat w}=\frac1{\sqrt{|\mathcal G|}}
    \sum_{z\in\mathcal G}\chi_w(z)\ket z,
 \qquad w\in\mathcal G,
\]
for the first output register of either field map. A row of the resulting
matrix is indexed by a frequency $w$ and a port $a\in\{1,\ldots,k\}$;
its columns are the $d$ input coordinates.

\begin{proposition}\label{fields:sparsity}
In these Fourier coordinates, every row of $\mathcal E_L$ and of
$\mathcal R_L$ has at most $L$ nonzero entries.
\end{proposition}

\begin{proof}
First, the sum of the vectors $v_i$ over any multiset of at most $L$
internal indices determines that multiset uniquely. Indeed, every coordinate
of such a sum, regarded as an ordinary nonnegative integer, is at most
$L2^{nL}<2^b$, so reduction modulo $2^b$ loses no information. Its first
coordinate records the length $m$, and the next $m$ coordinates record the
power sums $p_j=\sum_{t=1}^m i_t^j$. The elementary symmetric polynomials
are recovered successively from Newton's identities
\[
 e_0=1,\qquad
 j e_j=\sum_{t=1}^j(-1)^{t-1}e_{j-t}p_t\quad(1\le j\le m).
\]
These identities follow by taking the logarithmic derivative of
$\prod_{t=1}^m(1+i_t x)=\sum_{j=0}^m e_jx^j$ and comparing coefficients.
Consequently the monic polynomial
$T^m-e_1T^{m-1}+\cdots+(-1)^me_m$, and hence its multiset of roots,
is determined.

An input port contributes only to frequency zero and the same output port.
For an internal input $i$, the lower columns of $E_L^*$ are sums of
\[
 C^*Z^*,\quad C^*Z^*D^*Z^*,\quad\ldots,
\]
whereas those of $R_L^*$ are sums of
\[
 BZ,\quad BZDZ,\quad BZDZDZ,\quad\ldots.
\]
In both cases the rightmost diagonal factor acts on input $i$. Thus every
monomial contains the phase belonging to $i$, and its frequency is the
signed sum of the $v_j$ over a multiset of between $1$ and $L$ internal
indices containing $i$. The sign is negative for $\mathcal E_L$ and
positive for $\mathcal R_L$. For example, the second displayed term on
the $R$ side has entry
\[
 (BZDZ)_{ai}=\sum_j B_{aj}D_{ji}\,
                  \chi_{v_j+v_i}(z).
\]

Fix a nonzero frequency and an output port. For $\mathcal E_L$, negate
the frequency first to recover the corresponding positive tuple sum.
By the uniqueness just proved,
all its contributing monomials have the same multiset of internal indices.
Only indices belonging to that multiset can therefore occur as nonzero
input columns. There are at most $L$ of them. The ordering of indices in
a monomial can change its coefficient but cannot enlarge this column set.
No internal monomial has frequency zero, since its first coordinate has
absolute value between $1$ and $L$ modulo $2^b$. At frequency zero the
single port column gives the required bound as well.
\end{proof}

The role of the power-sum phases is now explicit: a frequency remembers
an unordered list of internal visits, so it permits only a short list of
possible input indices. The next section uses exactly this list to erase
the input register coherently.
